% Irodov Part III §3.5 251--262, user PDF pp.137--139,
% printed pp.128--130; figures 3.68--3.73 visually reviewed.
\electromagproblem{251}{128}
图3.68中细导线通电流 $I=8.0\,\mathrm A$。
求图示点 $O$ 处单位长度导线所受磁力模与方向：
(a) 半圆弯曲半径 $R=10\,\mathrm{cm}$；
(b) U 形两长平行段距离 $\ell=20\,\mathrm{cm}$。
\exsolution
局部导线受力每单位长度
$\bm f=I\,\bm t\times\bm B_{\rm others}$。
(a) 半圆弧在其中心产生
$B_{\rm arc}=\mu_0I/(4R)$，
直径段中心处的其他平直部分磁场不贡献，
故
\[
\boxed{f_a=\frac{\mu_0I^2}{4R}
\simeq0.20\,\mathrm{mN/m}}.
\]
(b) 图中点 $O$ 在左竖直连接段上，
每条半无限水平长直线距离它 $\ell/2$，
两条电流的磁场方向在该位置相加：
$B=2\mu_0I/[4\pi(\ell/2)]
=\mu_0I/(\pi\ell)$，
所以
\[
\boxed{f_b=\frac{\mu_0I^2}{\pi\ell}
\simeq0.13\,\mathrm{mN/m}}.
\]
\excheck 对每段导线应用
$I\mathrm d\bm\ell\times\bm B$，
按图中的电流箭头确定力朝向弯曲内部或外部；
图示 $O$ 不等于整个回路的质心。

\electromagproblem{252}{128}
单层紧密铜线线圈半径 $R=30\,\mathrm{mm}$，
导线直径 $d=0.10\,\mathrm{mm}$，
电流 $I=10\,\mathrm{mA}$，
线圈轴与均匀外磁场方向重合。
估算使绕组拉断的外部磁感应强度。
\exsolution
圆环上的磁力每单位导线长度为
$f=IB$，近似沿径向向外。
圆环微小弧元两端张力的径向合力
$T\,\mathrm d\theta$
与外力
$IBR\,\mathrm d\theta$ 平衡，
所以张力 $T=IBR$。
若铜线的极限抗拉应力为 $\sigma_m$，
允许最大张力
$T_{\max}=\sigma_m A=
\sigma_m\pi d^2/4$。
断裂条件：
\[
\boxed{B_{\rm rupture}=
\frac{\pi d^2\sigma_m}{4RI}}.
\]
以铜线强度数量级
$\sigma_m\sim3\times10^8\,\mathrm{Pa}$
估算：
$\boxed{B_{\rm rupture}\sim8\times10^3\,\mathrm T
=8\,\mathrm{kT}}$。
\excheck 题面未给具体铜材的抗拉强度，
因此数值必须注明采用的材料强度假设；
理想估算的8 kT远超普通实验室磁场。

\electromagproblem{253}{128}
截面积 $S=2.5\,\mathrm{mm^2}$
的铜线弯成正方形三条边，
可绕水平轴 OO' 转动
（图3.69）。外磁场竖直，
通电流 $I=16\,\mathrm A$ 后偏转
$\theta=20^\circ$。
求磁场模。
\exsolution
设正方形边长为 $a$，
铜密度为 $\rho_{\rm Cu}$。
图示三条边的重力矩由重心几何求得：
\[
\tau_g=2\rho_{\rm Cu}Sa^2g\sin\theta.
\]
电流在底部与两侧形成的有效磁力矩为
$\tau_B=Ia^2B\cos\theta$。
平衡：
$Ia^2B\cos\theta
=2\rho_{\rm Cu}Sa^2g\sin\theta$。
因此
\[
\boxed{B=\frac{2\rho_{\rm Cu}gS}{I}\tan\theta}.
\]
取 $\rho_{\rm Cu}\simeq8.9\times10^3
\,\mathrm{kg/m^3}$，
\[
\boxed{B\simeq1.0\times10^{-2}\,\mathrm T
=10\,\mathrm{mT}}.
\]
\excheck 两边含 $a^2$ 的矩相消，
故题目不需给正方形边长。

\electromagproblem{254}{128--129}
图3.70：$N=200$ 匝小线圈
位于磁极间，面积
$S=1.0\,\mathrm{cm^2}$，
以一端悬于天平横杆上，
力臂 $\ell=30\,\mathrm{cm}$。
线圈电流
$I=22\,\mathrm{mA}$。
追加配重
$\Delta m=60\,\mathrm{mg}$
才恢复平衡。
求该处磁感应强度。
\exsolution
线圈磁矩
$p_m=NIS$。
外磁场对它产生力矩模
$\tau_B=NISB$（图中转角使磁矩与场方向垂直）。
追加配重产生的重力矩
$\tau_g=\Delta mg\ell$。
平衡条件
$NISB=\Delta mg\ell$，
于是
\[
\boxed{B=\frac{\Delta mg\ell}{NIS}}.
\]
SI 单位：
$\Delta m=6.0\times10^{-5}\,\mathrm{kg}$，
$S=10^{-4}\,\mathrm{m^2}$，
\[
\boxed{B\simeq0.40\,\mathrm T}.
\]
\excheck 重量单位是毫克 mg，
不是克 g；换算错误会放大1000倍。

\electromagproblem{255}{129}
边长 $a=8.0\,\mathrm{cm}$ 的正方形电流框
$I=0.90\,\mathrm A$，
与无限长直导线电流
$I_0=5.0\,\mathrm A$ 同处平面。
电流框的转轴与长导线平行，
轴距是 $\eta a$，$\eta=1.5$。
求 (a) 合安培力；(b) 恒定电流下
绕该轴翻转 $180^\circ$ 的外功。
\exsolution
长线在距线 $r$ 处
$B(r)=\mu_0I_0/(2\pi r)$。
电流框的两条长边距直线分别为
$r_1=a(\eta-\tfrac12)$、
$r_2=a(\eta+\tfrac12)$。
两边受力方向相反，
其差为
\[
F=Ia[B(r_1)-B(r_2)]
=\frac{\mu_0II_0a}{2\pi}
\left(\frac1{r_1}-\frac1{r_2}\right),
\]
整理：
\[
\boxed{F=\frac{2\mu_0II_0}{\pi(4\eta^2-1)}
\simeq0.4\,\mathrm{\mu N}}.
\]
(b) 在恒定电流源约束下
两回路磁相互作用势的改变等于
需要的机械外功。
由电流框两长边位置积分长线磁场的磁通，
原平面构型与翻面构型的互感符号相反，
得到
\[
\boxed{A=\frac{\mu_0aII_0}{\pi}
\ln\frac{2\eta+1}{2\eta-1}
\simeq0.10\,\mathrm{\mu J}}.
\]
\excheck 本题要求机械外功而且\emph{电流恒定}；
不能把固定电流源提供的能量交换忽略。

\electromagproblem{256}{129}
两根很长平行导线半径为 $r$，
轴间距是 $\eta r$，$\eta=20$，
两导线电阻可忽略，
一端由外电阻 $R$ 连接，
另一端接直流电压源。
求两导线之间电磁合力为零时的 $R$。
\exsolution
两线带反号线电荷 $\pm\lambda$，
其横向\emph{静电吸引力}每单位长的模为
$F_e/L=\lambda^2/(2\pi\varepsilon_0d)$，
其中 $d=\eta r$。
两线由远端的电阻与电源闭合成回路，
故两长平行导线中的电流\emph{反向}，
其磁作用为\emph{排斥力}，
与电吸引力相反。磁力模
$F_m/L=\mu_0I^2/(2\pi d)$。
合力为零意味着
$\lambda/I=\sqrt{\mu_0\varepsilon_0}$。
两线每单位长电容
$C'\simeq\pi\varepsilon_0/\ln\eta$，
电压 $V=IR$，
$\lambda=C'V=C'IR$。
于是
\[
\boxed{R=
\frac{\sqrt{\mu_0\varepsilon_0}}{C'}
\simeq\frac{\ln\eta}{\pi}
\sqrt{\frac{\mu_0}{\varepsilon_0}}}.
\]
代 $\sqrt{\mu_0/\varepsilon_0}
\simeq377\,\Omega$，
\[
\boxed{R\simeq0.36\,\mathrm{k\Omega}}.
\]
\excheck 闭合回路使两线电流反向而磁斥，
同时两线所带静态线电荷相反而电吸，
二者才可能通过调节电阻相消。
若电源连接改变两作用方向，结论不再适用。

\electromagproblem{257}{129}
图3.61中的薄半圆截面长导体
通纵向电流 $I$，
另有中心细直线导体
沿\emph{相反方向}通相同电流。
求两导体间的单位长度磁相互作用力。
\exsolution
第225题已得到半圆电流截面
在中心的磁场模
$B_O=\mu_0I/(\pi^2R)$。
中心细线单位长度的磁力
$f=I B_O$，故
\[
\boxed{\frac FL=\frac{\mu_0I^2}{\pi^2R}}.
\]
\excheck 力方向沿半圆形截面对称轴，
由电流方向确定是朝半圆弧还是背离弧；
此处求的是模。

\electromagproblem{258}{129}
图3.71：两根无限长平行导体
通电流 $I_1,I_2$，
第一根为细线，第二根为宽 $b$
的均匀电流带，
最近距离为 $a$。
求两者单位长度磁相互作用力。
\exsolution
把宽带从近边缘 $x=a$
到远边缘 $x=a+b$
分割为并联细电流线，
每个宽度 $\mathrm dx$
中的电流
$\mathrm dI_2=(I_2/b)\mathrm dx$。
细线与其中一段的单位长度磁力
$\mathrm df/L=
\mu_0I_1\,\mathrm dI_2/(2\pi x)$。
积分：
\[
\frac FL=
\frac{\mu_0I_1I_2}{2\pi b}
\int_a^{a+b}\frac{\mathrm dx}{x}
=\boxed{\frac{\mu_0I_1I_2}{2\pi b}
\ln(1+b/a)}.
\]
\excheck 若 $b\ll a$，
对数约为 $b/a$，
恢复两条相距 $a$ 的细线公式。

\electromagproblem{259}{129}
两张无限平行电流片
形成仅在两片之间存在的均匀磁场 $B$，
片外磁场为零。求每片所受单位面积磁力。
\exsolution
对某一片，磁场的法向 Maxwell
应力差与两侧磁能密度差有关。
磁能密度 $w_m=B^2/(2\mu_0)$。
片内侧磁场为 $B$、片外为0，
故片面所受磁压力模
\[
\boxed{p_m=\frac FA=\frac{B^2}{2\mu_0}}.
\]
也可先求每片表面电流密度
$K=B/\mu_0$，
再取\emph{另一片}产生的局部场 $B/2$，
得 $p_m=K(B/2)$。
\excheck 计算自受力时不可直接
用本片两侧磁场的总值作用自身，
否则会多一倍。

\electromagproblem{260}{129--130}
图3.72(a)、(b)、(c)：
无限大电流片放在外加均匀磁场中，
其两侧总磁感应强度分别为
$B_1,B_2$。
求单位面积磁力及图示电流方向。
\exsolution
片面左、右两侧的磁压力分别为
$B_1^2/(2\mu_0)$、
$B_2^2/(2\mu_0)$。
因此沿从左侧指向右侧的正向，
单位面积合力为
\[
\boxed{\frac{F_x}{A}=
\frac{B_1^2-B_2^2}{2\mu_0}}.
\]
若图示 $B_1>B_2$，合力朝右。
电流方向由磁场切向跳变
$\bm K=\bm n\times
(\bm B_{\rm right}-\bm B_{\rm left})/\mu_0$
推定；图3.72的三组箭头对应
\emph{指向纸内}的表面电流。
\excheck 方向须结合图3.72的
\emph{磁场箭头}，
不能只由场强正数 $B_1,B_2$
决定表面电流的有向符号。

\electromagproblem{261}{130}
图3.73电磁泵中，熔融金属
在管道截面积近似为边长
$a=2.0\,\mathrm{cm}$ 的方形区域流动。
跨管道通电流 $I=100\,\mathrm A$，
磁场 $B=0.10\,\mathrm T$
垂直电流和管轴。求泵产生的表压。
\exsolution
电流在磁场中的总安培力
$F=IBa$，沿管道轴向推动金属。
若管截面积为 $a^2$，
泵的表压增量由
$\Delta p\,a^2=F$ 得
\[
\boxed{\Delta p=\frac{IB}{a}}.
\]
代入：
\[
\boxed{\Delta p=
\frac{(100)(0.10)}{0.020}
=500\,\mathrm{Pa}=0.50\,\mathrm{kPa}}.
\]
\excheck 表压正号方向取沿电磁泵
实际推动金属的方向。

\electromagproblem{262}{130}
长薄壁圆柱导体半径为 $R$，
沿轴向通总电流 $I$。
求导体圆柱壁受到的磁压力。
\exsolution
无限长薄壁圆柱内磁场为零，
外部磁场
$B_{\rm out}=\mu_0I/(2\pi R)$。
圆柱壁两侧磁能密度分别为
0 与
$B_{\rm out}^2/(2\mu_0)$。
磁压力（箍缩向内）的模
\[
\boxed{p_m=\frac{B_{\rm out}^2}{2\mu_0}
=\frac{\mu_0I^2}{8\pi^2R^2}}.
\]
\excheck 电流方向不影响压力大小，
因为压力正比于 $I^2$；
磁压力沿径向指向圆柱轴。
